\section{Discussion}
\subsection{What the results say}
The decomposition reframes a familiar complaint. The literature on causal readings of attributions is usually stated as a warning: SHAP values are not causal effects. Our benchmark puts a number on the warning and locates it. The largest component of the gap between an explanation and the causal effect is not the explainer's approximation and not even the model's imperfection; it is the \emph{definition} of what an attribution measures. This has a practical consequence: investing in a better Shapley estimator (more coalitions, exact TreeSHAP instead of KernelSHAP) reduces the smallest term of~\eqref{eq:decomp} and leaves the dominant one untouched. Only a different estimand---causal or asymmetric Shapley values with a known graph~\cite{heskes2020,frye2020}---can act on it, and those require the very knowledge that is missing in practice.

The second result is more optimistic. Although the gap is structural, its \emph{size} correlates with properties that are visible without the graph. The observables that carry the most signal are of two kinds. Disagreement signals (between explainers on the same model, and between the two model families for the same explainer) are high exactly when the fitted response surface is ambiguous about which correlated inputs to credit, which is also when the credit assignment is most likely to depart from the interventional effect. Fit and concentration signals (test $R^2$, prediction spread, attribution Gini) indicate how much of the outcome the model captures and how many features it leans on; a model that spreads credit widely over correlated inputs has more room to misrank them. Multicollinearity by itself, the signal practitioners most often check, is comparatively weak on its own but useful in combination.

The diagnostic transfers across graph families because it is not learning the families: it has no access to the family label, the mediation share or the presence of confounding, and the LOGFO protocol removes the held-out family's structure from training entirely. What it learns is a mapping from symptoms to error that the five families share, which is the property a diagnostic needs if it is to be used on data whose causal structure is unknown.

\subsection{Limitations}
\textbf{Scope of the benchmark.} All SCMs are small ($d \le 10$, so that exact Shapley values can be enumerated), have additive Gaussian noise, a single continuous outcome, and come from five hand-designed families. Real tabular data has categorical features, heteroscedastic noise, selection effects and far more features. The diagnostic is calibrated to the benchmark's distribution of situations; it should be re-trained or at least re-validated before being trusted elsewhere, and the semi-synthetic checks we planned (IHDP, Adult Income, California Housing) remain future work.

\textbf{Choice of estimand.} We measure error against one causal quantity, the total effect of a unit-standard-deviation shift. Direct effects, binary treatments, or effects on a conditional sub-population would give different estimand gaps, and some would be closer to what interventional SHAP measures. The decomposition itself does not depend on this choice; the numbers do.

\textbf{What the diagnostic does and does not do.} It predicts how far an explanation is likely to be from the causal ranking; it does not say which features are misranked or what the correct ranking is. A high predicted error is a reason to seek domain knowledge or an experiment, not a correction. It is also a regression on the explanation's error magnitude; we report AUROC for a median split as a convenience for a warning threshold, not as a calibrated probability.

\textbf{Dependence between rows.} Rows that share an SCM share data and ground truth. All evaluation is grouped by SCM and all tests resample SCMs, but the effective sample size for the family-level conclusions is 60 SCMs per family, and the per-family LOGFO numbers should be read with that in mind.

\textbf{Compute.} Exact Shapley values of the fitted model cost $100 \times 50 \times 2^d$ model evaluations per seed; the full run (300 SCMs, 2 models, 3 seeds, 5 explainers, perturbation re-explanations) takes about three hours on two CPU cores. Larger $d$ would need sampled-coalition oracles, which would blur the line between $\Phi(f)$ and $\hat E$.
